\section{Discussion and Limitations}\label{sec:discussion}
The central question is how to adapt VLA-based control when an existing action interface does not provide useful behavior on a new robot. Dense reward supports skill discovery, binary refinement targets completion, self-imitation consolidates experience, and staged RL learns the complete instruction. These roles require separate measurements of acquisition, retention, and composition.

The studies meet at a narrower mechanism: ordered task achievements can distinguish trajectories that share the same terminal failure. CDPR exposes these distinctions through milestone advantages; LIBERO also resolves progress within a phase. Its observed dense advantage supports the value of intermediate feedback beyond initial acquisition. The choice of its granularity remains empirical: the LIBERO ablations do not isolate progress density, and CDPR lacks a matched dense-versus-milestone composition comparison.

There is therefore no requirement to remove all intermediate feedback once skills are learned. Binary skill refinement and milestone composition serve different objectives and reset distributions. The schedule uses binary rewards to target skill completion, then reintroduces structured feedback for the longer composed task. LIBERO makes that design rationale plausible; it does not establish that switching to binary rewards, or the full schedule, is superior to continued dense training.

The CDPR is one embodiment-adaptation case, using a new residual over a retained SmolVLA expert. Resetting an entire expert or applying the full schedule to additional embodiments remains untested; LIBERO starts with a trained policy. Skill consolidation also gives no guarantee that further SFT on full-task successes helps: unsuccessful campaign attempts show the importance of data mixture, coverage, and action consistency.

The current simulation endpoint uses development scenes and one training lineage. Matched schedule, self-imitation, and composition ablations, independent seeds, and a locked final test remain necessary for stronger attribution. The LIBERO draft lacks trial denominators and quantitative throughput records. Physical success remains the author's drafting assumption until completed trial evidence is inserted. Manual rewards, historical scripted initialization, calibration, and inherited training costs also remain part of the method's scope.
